\begin{problem}[A mixed arithmetic-geometric local ring]
Let $A=\ZZ[x]$ and $\mathfrak m=(2,x-1)$.  Prove $\mathfrak m$ is maximal, compute $\kappa(\mathfrak m)$, characterize units of $A_{\mathfrak m}$, and compare this point with the generic prime $(0)$ and the vertical prime $(2)$.
\end{problem}

\paragraph{Why this problem matters.}
It combines arithmetic and geometric localization in one example and mirrors the legacy mixed-ring computations.

\paragraph{Useful ideas before starting.}
Reduce coefficients modulo $2$ and set $x=1$.

\paragraph{Guided hints.}
The quotient by $\mathfrak m$ should be a field of two elements.

\begin{solution}
The evaluation-reduction map $\ZZ[x]\to\FF_2$ sending coefficients mod $2$ and $x\mapsto1$ is surjective with kernel $(2,x-1)$, so $\mathfrak m$ is maximal and $\kappa(\mathfrak m)=\FF_2$.  In $A_{\mathfrak m}$, an element is a unit exactly when it is not in $\mathfrak m$, equivalently its image in $\FF_2$ is nonzero.  The generic point $(0)$ has local ring $\QQ(x)$, while $(2)$ has residue field $\FF_2(x)$ and represents the generic point of the special fiber over $2$.
\end{solution}

\paragraph{What happened mathematically.}
The same spectrum contains a generic characteristic-zero point, a characteristic-$2$ generic fiber point, and a closed characteristic-$2$ point.

\paragraph{Extensions.}
Replace $2$ by a prime $p$ and $x-1$ by an irreducible polynomial modulo $p$.

\paragraph{Provenance.}
Expands the mixed $\ZZ[x]$ local-ring example in legacy \emph{Theory of Commutative Algebra 13}.
